La separación de fuentes musicales consiste en descomponer una mezcla musical en varias señales individuales, por ejemplo, voz, batería, bajo y otros.

Una formulación clásica del problema es la separación en dos fuentes, vocal y acompañamiento, mientras que herramientas más modernas permiten extenderse a pistas mas complejas y devolver más elementos.
Las dificultades clásicas del problema se pueden resumir así\cite{rafii18}:

- Solapamiento espectral: las fuentes suelen compartir frecuencias, especialmente cuando existen armónicos extendidos por todo el campo espectral
- Fase: la información de fase es crucial para la calidad de la separación, pero es difícil de modelar y recuperar.
- Mezcla comercial: las mezclas comerciales suelen tener efectos de reverberación, compresión y otros procesos que complican la separación.
- Simultaneidad de instrumentos: en una mezcla musical, normalmente, los instrumentos suenan simultáneamente, lo que dificulta la identificación y separación de cada fuente.

